The column that came within sight of King’s Landing’s walls, on the forty-first day since Winterfell fell, was not an army any longer, not in the sense the word had meant when Tyrion Lannister first watched it march north from Dragonstone with three dragons and ten thousand disciplined spears and a queen who still believed, then, that conquest and liberation were the same thing worn in different clothes.

He stood at the rise above the kingsroad’s final descent and made himself count properly, because he had learned, across a long career of managing other men’s wars, that a Hand who flinched from the arithmetic was a Hand who deserved to lose. The Unsullied numbered perhaps three thousand of the seven that had sailed from Essos, their ranks thinned by Winterfell and by three subsequent rearguard actions and by a slow, grinding sickness that had taken more of them these past weeks than any battle had. The Dothraki numbered fewer still — bloodriders on foot now as often as mounted, their horses eaten or lost or simply worked to death along a road that had never once been kind to creatures bred for open grass and warm wind. The Northern levies had lost their lords in twos and threes the whole retreat through, Robett Glover among the first, and the smallfolk who had marched rather than ridden numbered so few, compared to the column that had left Wintertown, that Tyrion had stopped asking Maester Wolkan for the precise count some while back, unable to bear hearing it said aloud each week smaller than the week before.

One dragon. Drogon, circling now above the column in the cold morning light, the last living child of a queen who had crossed an ocean with three. Tyrion had watched Daenerys track him with her eyes more often than she tracked anything else these past weeks, the way a woman watches the one thing left of a family that had once, briefly, numbered considerably more.

They had outrun the dead. That much, at least, the arithmetic still favored them — Melisandre’s fire and Jon’s two hundred and a dozen smaller sacrifices since had bought them a lead that held, though every scout’s report these past days had grown steadily more urgent about how much of that lead remained. And ahead of them now, rising grey and vast out of a grey winter sky, sat the city Tyrion had been born to rule someday, in whatever fashion the House of Lannister had once imagined its second son might rule anything at all.

\sceneBreak

They sent the heralds forward at midday, under a peace banner that Tyrion privately doubted would matter to a woman who had burned a Sept full of septons and smallfolk alike rather than answer for her own crimes before a court she didn’t control. Jon rode with them, and Tyrion rode beside him, because he understood, better than anyone else in that column, exactly what shape his sister’s welcome was likely to take, and found he could not let Jon face it without at least the small comfort of a Lannister who already knew.

They stopped a bowshot from the great gate, close enough to see the walls properly — manned, thick with soldiers Tyrion recognized by their gear as Golden Company rather than the city’s own thin garrison, and scattered among them, at intervals he did not like the careful spacing of, the dark iron shapes of scorpions, the same design that had once been used to bring dragons down over the Reach.

“They know we’re weak,” Tyrion said, low, to Jon, reading the placement of those weapons with a sinking certainty that had nothing to do with guesswork. “Those scorpions aren’t spread to defend a city from siege engines, Jon. They’re spread to bring down a single dragon flying reconnaissance. Someone’s told them exactly how thin we’ve worn, and told them recently enough that they’ve had time to prepare for it properly.”

Jon said nothing, though something tightened in his jaw that told Tyrion he’d reached the same conclusion by a different road.

The gate did not open. Instead, after a long wait that Tyrion understood was itself a message — a queen making an army stand in the cold to demonstrate whose patience mattered less — a small side postern opened, and a herald came out alone, richly dressed against the weather in a way that only sharpened, by contrast, how threadbare the column behind Tyrion had become.

“Her Grace, Queen Cersei of House Lannister, First of Her Name,” the herald announced, in the flat, over-practiced tone of a man reciting words he’d been made to memorize rather than words he believed, “bids the Dragon Queen and her allies welcome to the gates of King’s Landing, and regrets that circumstances do not presently permit her to open them.”

\sceneBreak

“Circumstances,” Tyrion repeated, before Jon could speak, because he understood his sister’s particular cruelties better than anyone else present and recognized, in that single bloodless word, the shape of exactly how much damage she intended to do with it. “Tell my sister I’d hear the circumstances explained properly, rather than dressed in a herald’s borrowed courtesy.”

The herald’s eyes flicked to him — recognition, plainly, though no particular warmth in it. “Her Grace bids me say this, my lord, word for word, as she gave it to me. That she has watched a column of the dead march south from the Wall for near two months now, gaining on whatever remains of the army presently standing at her gates. That she does not believe an army so reduced — by her own counting, less than half the strength that marched north, one dragon of three, and starving besides — stands any real chance of defeating what pursues it. And that a queen who has already buried three children does not intend to gamble a fourth, and a city of half a million souls beside, on opening her gates to host a battle she believes is already lost before it’s fought.”

Tyrion felt the words land with a precision that told him, beyond any remaining doubt, exactly where Cersei’s numbers had come from. \textit{By her own counting.} No guess would have produced figures that close to true. Someone in this column, or someone who had recently left it, had given his sister an accounting accurate enough to build a wall’s worth of scorpions around.

“She means to let us fight the dead alone,” Jon said, flat, “and take whatever’s left standing afterward.”

“She means precisely that,” Tyrion said, “and I confess, examining the strategy purely as strategy and setting aside for one difficult moment that the woman proposing it is my own sister, I find I cannot fault the cold arithmetic of it. Two enemies march toward her walls. She need spend nothing on either, if she simply waits, and take whatever’s left of both once they’ve bled each other properly.” He heard his own voice go rougher than he intended. “It is, in fact, precisely the sort of strategy our father would have praised her for, which I imagine is exactly why she chose it.”

\sceneBreak

The herald had more to say, and Tyrion understood, watching his careful, rehearsed face, that the worst of it was still to come.

“Her Grace further bids me say,” the herald continued, “that any lord presently marching with this column who wishes to be received individually, and separately, within the city’s walls — bringing with him such men and such intelligence as remain useful to Her Grace’s cause — will find a considerably warmer welcome than the column as a whole is presently offered.”

The silence that followed told Tyrion that every man within hearing understood exactly what was being offered, and to whom, and Tyrion felt something cold settle into his stomach that had nothing to do with the winter air — because a queen did not extend an invitation like that one unless she already had reason to believe at least one lord in this column might accept it, and Tyrion understood, in the same breath, that the intelligence behind her walls and the temptation now being dangled before them both pointed toward the same unhappy conclusion: someone in this column had already sold them, and Cersei was simply inviting him to formalize the arrangement.

“Tell your queen,” Jon said, before Tyrion could compose a properly diplomatic answer, “that this column did not march six hundred miles through starvation and grief to be picked apart by an invitation to treachery. Tell her we’ll camp at her gates, and hold what ground we can, and that when the dead arrive — and they will arrive, whatever her scorpions are aimed at — she’ll discover exactly how thin the wall between her safety and their numbers actually is, whether she’s opened her gates to help fight them or not.”

The herald bowed, unmoved, a man long past being troubled by the defiance of the doomed, and withdrew back through the postern, and the great gate remained shut, and Tyrion stood in the cold with Jon beside him and felt, for the first time since Winterfell fell, something closer to despair than he had permitted himself across the whole grinding length of the retreat.

\sceneBreak

The war council gathered that evening in a tent pitched within sight of the walls that would not open to them, and Tyrion watched the faces around the fire — Jon’s controlled fury, Sansa’s flat, calculating stillness, Daenerys’s careful, dangerous quiet, Varys watching everyone with his particular unreadable attention — and understood that they had marched out of one trap only to walk directly into a second, more patient one.

“We can’t besiege a city with the dead a fortnight behind us,” Davos said, laying the problem out plain, the way he always did. “And we can’t outrun the dead by marching further south, not with the state our own numbers are in. We’re caught, plain as that. Dead behind, shut gates ahead, and a queen who’s made it clear she means to let the two of us bleed each other white while she watches from behind stone.”

“Then we don’t besiege her walls,” Sansa said, quiet. “We use them. Whatever else Cersei intends, those walls are the strongest fortification in Westeros, and the dead will have to break against them the same as they broke against Winterfell’s, whether she opens her gates to us or not. If we position properly — dig in at her doorstep, force her hand by making the battle happen whether she wills it or not — she may find her cold arithmetic considerably harder to hold to once the dead are actually climbing her own walls rather than merely camped conveniently outside them.”

“She’ll have accounted for that,” Tyrion said, though he found, saying it, that some small, cornered part of him was already reaching for exactly the plan Sansa had just proposed, because he could not presently think of a better one. “My sister does not build a strategy without accounting for the obvious counter to it. But I agree we’ve no better option before us. We dig in. We hold what ground we can between her walls and whatever comes down that kingsroad after us. And we make it as costly as possible for her to go on pretending this isn’t her war too, once the dead are close enough that her own scorpions have to choose a target.”

Daenerys had said little through the whole council, and Tyrion watched her now, in the firelight, some old fierce calculation moving behind her eyes that he had learned, these past weeks, not to interrupt until she was ready to share it.

“She’s made her choice,” Daenerys said finally. “Let her live with it. We hold this ground, we survive what’s coming however we manage it, and when this is finished — if any of us are left to finish it — I intend to remind Cersei Lannister precisely what her cold arithmetic cost the people who trusted her to at least pretend to be part of this war.” Something in her voice had gone very quiet, very still, the particular quality Tyrion had come to recognize, this past month, as more dangerous than any shout. “I’ve buried a great deal these past weeks, or watched it turned into something worse than dead. I find I’ve very little patience left for queens who believe watching from behind a wall makes them exempt from what the rest of us are paying to survive.”

No one answered that, because there was nothing left to say that would not simply confirm what every face around that fire already understood — that they had reached the end of the road they could retreat along, and the beginning of one they would now have to hold, whatever waited behind them and whatever refused to help them from the walls ahead.

\sceneBreak

Later, alone, Tyrion walked the edge of the camp where it faced the city, close enough to see torchlight moving along the battlements, and found himself doing the thing he had spent the whole council avoiding — thinking of Cersei not as a strategist to be countered but as his sister, the girl who’d once, in some other lifetime neither of them could properly remember without wincing, actually loved him before their father taught her better.

He thought of the herald’s offer, the invitation to treachery dangled before the whole column, and understood, turning it over, that it was not merely a tactic. It was a message meant for him specifically, whatever else it accomplished — a reminder, from a sister who knew him better than anyone living, that she still believed, after everything, that every man had his price, and that a dwarf who’d spent his whole life being told he was worth less than his golden siblings might yet be tempted to prove her right one final time.

He was not tempted. He understood that plainly, standing in the cold, and found the understanding did not comfort him nearly as much as he’d have expected. Because the not-being-tempted meant he had finally, fully let go of whatever remained of the boy who’d once wanted nothing more than his sister’s approval, and some part of him — smaller than he liked, more honest than he liked — mourned that letting go the way a man mourns a wound he knows was necessary to close.

He thought of Podrick, gone. Of two hundred men who’d held a narrows and were gone. Of a whole column of smallfolk who’d marched out of Wintertown and did not number, now, even a third of what had left. He thought of Cersei behind her walls, warm, fed, safe, watching two armies she believed doomed regardless approach each other across her doorstep, confident in an arithmetic that had never once, in her whole life, accounted properly for the cost of being wrong.

\textit{Someone in this column sold us.} He turned that over too, last of all, because it was the thought he liked least and therefore the one he trusted himself to examine most carefully. He did not yet know who. He suspected, from the particular precision of the numbers the herald had recited, that it was a lord rather than a common soldier — someone with access to the council’s own figures, someone hungry or bitter or frightened enough that Cersei’s gold, or Cersei’s promises, had outweighed whatever remained of his loyalty to a cause that had cost him everything already and promised him nothing certain in return.

He did not know it was Rowan Glover. He would learn that later, and would remember, learning it, this exact night, standing at the edge of a camp that could not go forward and could not go back, wondering which of the exhausted, grieving, starving faces around that dying fire had already decided the war was lost and chosen, quietly, to be on the winning side of its losing.

Above him, Drogon circled once more against the darkening sky, alone, the last of three, and Tyrion Lannister turned back toward the camp at last, toward a war council that had run out of roads to retreat along, and understood that whatever came next — siege, battle, betrayal fully unmasked, or the dead finally catching them where they stood — it was going to happen here, at his sister’s shut gates, whether she ever opened them or not.